\section{Introduction}
\label{sec:introduction}

Reasoning has become a central axis of large language model development. With each generation, models train on more and better data with improved training methods, and become capable of tackling problems that would have been out of reach a few years ago. Their internals, in turn, grow more informative: a richer internal computation is doing the work behind richer external behavior, which makes the question of what those internals encode increasingly worth asking. Code generation is a particularly clean field for that question: outputs are easily verifiable against unit tests, with no human-judgment loop required to decide whether an answer is correct. As model performance on code-generation benchmarks has improved rapidly, the question is no longer only \emph{whether} a model can produce correct code, but also \emph{what happens internally} as it does so. This shifts the focus toward mechanistic interpretability in models that now exhibit genuine capability in code generation.

A standard tool for asking what is encoded inside a frozen language model is the \emph{linear probe}: a simple linear classifier trained on a hidden state vector to predict some property of interest, where probe performance is read as a measure of how linearly decodable that property is from the representation \citep{alain2016probes}. Two recent works have studied code correctness through hidden states \citep{bui2025openia, ribeiro2025internal}, reading hidden states of already-generated code to judge whether it is correct without running the unit tests. One trains a classifier on those states; the other contrasts correct and incorrect samples to extract a correctness direction and rank candidate generations. In both, the hidden state is read \emph{after} the model has produced an answer.

This work asks a related but earlier question: whether correctness is already decodable from the hidden state at the position of the last prompt token, captured on a single forward pass over the prompt and before any output token is sampled. Whereas those works operate on hidden states of generated code, the probe in this paper operates on the hidden state of the prompt alone. This is a more upstream interpretability claim: the information must already be present at the moment the model finishes reading the input, not just at the moment it has committed to a particular output.

A second, more geometric question follows. When the model is given the chance to repair its own failing code, does the shift in its hidden state between the failed attempt and the repair attempt carry a stable signature of whether the repair will succeed? Phrased this way, the question lives within the Representation Engineering framework \citep{zou2023repe}, which identifies candidate features as \emph{directions} in activation space, typically as the difference between two group means. The natural application here is to compute the per-task hidden state delta from a failing attempt to its repair attempt, split the deltas by repair outcome, and ask whether the contrastive direction is stable.

Such contrastive constructions, however, are subtler than they look. A direction of the form $\mathrm{mean}(\text{success}) - \mathrm{mean}(\text{failure})$ already cancels any covariate whose mean is equal across the two groups, so residualizing the deltas against such a covariate would only collapse the signal artifactually. Conversely, when a covariate's group means differ, the raw direction is partially driven by it rather than by the latent feature of interest. The methodological approach of this paper makes this effect explicit: a short distribution test per candidate covariate, run before any residualization, decides which residualizations are warranted and which would be misleading.

The two findings are the following. The pre-generation probe holds up cleanly: the prompt-final hidden state carries a linearly decodable correctness signal that is robust across many random train/test splits and that survives controlling for prompt length. The repair direction is likewise statistically detectable in the raw deltas, passing both magnitude and inclination tests against label-shuffled nulls, and it remains significant after a conditional residualization once three observable repair-context covariates are checked and found to differ between the success and failure groups, even as that residualization removes a substantial part of it. The honest reading is that the geometric direction is a genuine repair-success signal that the observable repair context does not fully account for, rather than a mere artifact of that context. Consistent with prior reports that iterative self-repair plateaus after the first one or two attempts \citep{arimbur2026howmanytries}, the per-attempt pass rate collapses sharply after the first repair, and later transitions are too data-starved to support the same direction analysis.

Read together, the two results illustrate how the same confound control procedure applies to two different signals. For the pre-generation correctness probe, residualizing against prompt length confirms that the decodable signal is not reducible to prompt structure, even if a small part of the raw probe's performance is. For the geometric repair direction, residualizing against the differentially distributed repair-context covariates removes the part of the signal that the repair context can account for, and a repair-success signal remains once that part is taken out. In both cases the procedure measures how much of a raw signal is genuinely its own.